% First complete draft -- China's out-of-home food consumption
% Assembled by lead editor from PhD1--PhD4 sections. Target ~9000 words.
% For: American Journal of Agricultural Economics, Agricultural Economics, Journal of Agricultural Economics.

\documentclass[12pt,a4paper]{article}
\usepackage[utf8]{inputenc}
\usepackage[T1]{fontenc}
\usepackage{amsmath,amssymb}
\usepackage{booktabs}
\usepackage{natbib}
\usepackage{geometry}
\usepackage{hyperref}
\usepackage{parskip}
\geometry{margin=2.5cm}

\title{Machine-Learning-Based Estimates of Out-of-Home Food Consumption in China: Provincial and National Coefficients for Food Security Assessment}
\author{[Authors to be inserted]}
\date{\today}

\begin{document}
\maketitle

\begin{abstract}
Existing work on China's food away from home (FAH) or out-of-home food consumption has relied mainly on household surveys and expenditure data, documenting rising FAH shares of total food spending and income-driven growth in urban dining-out, but has not produced nationally or provincially representative estimates of out-of-home consumption by food commodity that can be used for macro-level food balance and food security analysis. International sources such as FAO, OECD, and IFPRI report aggregate food use or consumption trends for China yet typically do not distinguish at-home and out-of-home consumption by product. This paper fills that gap by providing the first machine-learning-based, commodity-level estimates of out-of-home consumption coefficients for China at provincial and national levels. We combine micro household survey data with macro province-year covariates through Copula income-distribution matching, imputation of missing consumption ratios, spatial interpolation for provinces without micro data, and a comparison of 12 machine-learning models with bootstrap uncertainty quantification. The primary deliverable is coefficient estimates (total-to-at-home ratio by category and province-year), from which per-capita consumption quantities are derived in post-processing and reported in the results. We report national and provincial out-of-home consumption coefficients for 15 food categories over 2015--2024 and bootstrap-based 95\% prediction intervals (reported in the results and supplementary material). Main results show coefficients above one for all categories, with the largest relative contributions from poultry, pork, legumes, and coarse grains. Findings are robust to excluding imputation and support the use of total (at-home plus out-of-home) consumption in grain and food demand assessments and in discussions of China's food security.
\end{abstract}

\noindent\textbf{Keywords:} food away from home; China; food consumption; food security; machine learning; Copula; spatial interpolation; bootstrap.

% ========== 1. Introduction ==========
\section{Introduction}
\label{sec:introduction}

Accurate measurement of how much food is consumed away from home is essential for understanding dietary change, assessing food security, and designing effective agricultural and food policy. In China, the share of food expenditure devoted to eating out has risen sharply over recent decades, and out-of-home consumption now constitutes a substantial part of total food demand. Yet existing studies have not delivered nationally or provincially representative, commodity-level out-of-home consumption coefficients that can be used in macro-level food balance sheets and demand projections. International datasets (FAO, OECD, IFPRI) report aggregate consumption or use for China but do not separate at-home and out-of-home consumption by product. This gap limits the accuracy of food demand and food security assessments that rely on at-home consumption alone.

This paper addresses the gap by producing the first set of machine-learning-based, commodity-level out-of-home consumption coefficients for China at provincial and national level. We define the out-of-home consumption coefficient as the ratio of total (at-home plus out-of-home) consumption to at-home consumption, by food category and province-year. Using micro household survey data to train prediction models and Copula-based income distribution matching to bridge micro and macro data, we estimate coefficients for 15 food categories (four grain types, vegetables, meat, aquatic products, eggs, dairy, fruit, edible oil, sugar) for years 2015--2024. We compare 12 machine-learning models, select the best model for imputing missing consumption ratios, apply spatial interpolation for provinces without micro data, and report robustness checks (no imputation) and bootstrap prediction intervals. The resulting coefficients can be used directly in food balance and food security analysis and can be compared with FAO, OECD, IFPRI, and the survey-based literature where definitions overlap.

The main innovations are: (1) Copula income distribution matching to align micro and macro feature space so that models trained on household data can be applied to province-year aggregates; (2) systematic comparison of 12 ML models for coefficient prediction and for imputation, with the best imputation model (FT-Transformer) used to produce a unified imputed sample; (3) spatial interpolation (IDW for coefficients, with production shares for grain structure) for provinces missing in the micro data; (4) robustness to dropping imputation and bootstrap-based uncertainty quantification; and (5) the first ML-based, commodity-level out-of-home coefficients for China at provincial and national level.

The remainder of the paper is organized as follows. Section~\ref{sec:literature} reviews concepts and empirical evidence on food away from home, with emphasis on China, and states the gap this paper fills. Section~\ref{sec:research-design} sets out the research question, data and methodological difficulties, and the chosen research strategy. Section~\ref{sec:methods} describes the prediction target, features, the 12 ML models, and the pipeline (imputation selection, coefficient prediction, spatial interpolation, post-processing, robustness, and bootstrap). Section~\ref{sec:results} presents main estimates, imputation selection results, and robustness. Section~\ref{sec:discussion} discusses the estimates in light of FAO, OECD, IFPRI, and the literature, and their implications for China's food security. Section~\ref{sec:conclusion} concludes and Section~\ref{sec:policy} offers policy recommendations.

% ========== 2. Literature Review (PhD1) ==========
\section{Literature Review}
\label{sec:literature}

% Literature review section (博士生1) — based on lit folder
% Standalone section: can be \input after \section{Literature Review} in main file.
% Citations refer to references in paper_draft/references.bib (lit folder entries added).

\subsection{Concepts and Measurement of Food Away from Home}

Food away from home (FAH), or out-of-home food consumption, refers to food prepared and consumed outside the household. Standard definitions in agricultural and consumer economics include meals and snacks purchased at restaurants, canteens, and other out-of-home establishments, and sometimes extend to institutional and business-provided meals \citep{USDA_ERS_FAH}. Measurement can be in \emph{expenditure} (spending on FAH) or in \emph{quantity} (physical consumption of food items consumed away from home). Expenditure-based measures are widely available from household budget and consumer expenditure surveys and have been used in a long line of work on FAH demand. Early US studies examined household consumption of food-away-from-home by total expenditure and by type of facility \citep{McCracken1987}, participation and expenditure patterns \citep{Byrne1996}, and the role of working wives and demographic factors \citep{Kinsey1983, Yen1993}. Similar approaches have been applied in other countries, for example panel data analysis of demand for food away from home in Spain \citep{Angulo2007}. Quantity-based measures require survey instruments that record both home and total (or out-of-home) consumption by food category, which are less common and harder to aggregate to the macro level.

For policy and food security analysis, a key quantity-based indicator is the ratio of total consumption to at-home consumption (or its inverse, the at-home share). In this paper, the \emph{out-of-home consumption coefficient} is defined as total consumption divided by at-home consumption, so that total food use can be inferred from at-home consumption data when the latter is the only dimension available in macro or regional statistics. This coefficient is distinct from FAH expenditure shares and elasticities commonly reported in the FAH literature and is the quantity-based input needed for food balance and demand projections.

\subsection{Empirical Evidence on China's Food Away from Home}

Evidence on China's FAH consumption comes mainly from household surveys and expenditure studies. \citet{Ma2006} showed that consumption of food away from home in urban China rises with income (``getting rich and eating out'') and documented strong income gradients using urban household data. \citet{Gould2006} assessed the current structure of food demand in urban China and provided estimates of demand elasticities that inform understanding of how food consumption responds to income and prices. Disaggregated evidence by type of facility and meal is available for Beijing: \citet{Bai2012} disaggregated household expenditures on food away from home in Beijing by type of food facility and type of meal, showing heterogeneity in where and how urban households consume food away from home. \citet{Liu2015} examined the role of household composition and income in food-away-from-home expenditure in urban China, reinforcing that demographics and income are key drivers. More recently, \citet{Zheng2019} predicted the changes in the structure of food demand in China, including the role of FAH in evolving diets. These studies establish that China's FAH is substantial and growing, that income and urbanization are primary drivers, and that household composition and demographics matter. They rely largely on expenditure or share measures from urban household surveys and the China Health and Nutrition Survey (CHNS); they do not produce nationally or provincially representative, commodity-level \emph{quantity} coefficients (total-to-at-home ratios) suitable for macro-level food balance and food security modeling.

\subsection{International Evidence and Methodological Approaches}

International work on FAH demand provides useful methodological benchmarks and confirms that the distinction between food at home and food away from home matters for demand and policy analysis. \citet{Balagtas2023} analyze consumer demand for food at home and food away from home and discuss responsiveness to policy. \citet{Brauw2020} address income variability and evolving diets in elasticity estimation for processed foods. Studies in other settings, such as consumption of processed food and food away from home in big cities, small towns, and rural areas \citep{Sauer2021}, show that settlement type and geography affect FAH patterns and that comparable quantity-oriented measures are scarce. Shocks such as the COVID-19 pandemic have been shown to affect food purchase behavior and food values \citep{Ellison2020}, underscoring the need for robust and updatable coefficient estimates. Together, this body of work supports the importance of FAH in modern food systems but also highlights that commodity-level, quantity-based out-of-home coefficients are rarely available for use in national or regional food balance and food security assessment.

\subsection{Data and Methodological Limitations}

Three limitations are central. First, \emph{micro--macro mismatch}: household surveys contain individual or household income and detailed consumption, whereas macro or provincial data typically provide only average income and aggregate or at-home consumption. Using micro elasticities or coefficients directly for macro prediction therefore requires a way to map macro covariates to the distribution of behavior observed in the micro data. Second, \emph{spatial and product coverage}: many surveys do not cover all provinces or all years, so that provincial-by-commodity out-of-home coefficients are missing for some region--year cells and must be inferred or interpolated. Third, \emph{measurement and model choice}: different surveys define ``at home'' and ``away from home'' differently, and the choice of functional form or model can affect estimated elasticities and implied coefficients. These issues complicate the production of consistent, commodity-level, provincial or national FAH coefficients suitable for food balance and policy analysis.

\subsection{The Gap Addressed by This Paper}

Existing studies have established that China's FAH consumption is substantial and growing and have identified correlates (income, demographics, urbanization) and data sources (urban household surveys, CHNS, expenditure surveys). However, they have not produced nationally or provincially representative, commodity-level out-of-home consumption \emph{coefficients} (quantity-based, total-to-at-home ratio) suitable for macro-level food demand and food security modeling. No existing study or international dataset provides such coefficients for China. This paper addresses the gap by: (1) using micro survey data (CHNS) to train machine-learning models for the at-home to total consumption ratio (and thus the out-of-home coefficient) by food category; (2) bridging micro and macro data via income-distribution matching so that predictions can be made for province-year cells with only macro covariates; (3) applying spatial interpolation and, where relevant, production structure to handle missing provinces; (4) comparing multiple ML models and reporting bootstrap-based prediction intervals; and (5) producing the first set of ML-based, commodity-level out-of-home consumption coefficients for China at provincial and national levels for use in food balance and food security assessments.


% ========== 3. Research Design (PhD2) ==========
% Research design section (博士生2)
% Paper: How much does China consume food out-of-home?
% Target: Top agricultural economics SSCI journals

\section{Research Design}
\label{sec:research-design}

\subsection{Research Question}
\label{subsec:research-question}

The central research question of this paper is: \textbf{How much does China consume food out-of-home?} We operationalize this at two levels and by food category. First, at the \textbf{national} level, we aim to estimate aggregate out-of-home food consumption and its evolution over time. Second, at the \textbf{provincial} level, we aim to estimate out-of-home consumption by province and year, so as to capture spatial heterogeneity and support region-specific policy analysis. In both cases, estimates are produced by \textbf{food category}: four grain types (rice, wheat, beans, coarse grains), vegetables, meat (pork, beef, mutton, poultry), and other categories (aquatic products, eggs, dairy, fruit, edible oil, sugar)---15 categories in total. The key outcome is the \emph{out-of-home consumption coefficient} (ratio of total to at-home consumption), from which provincial and national at-home and out-of-home consumption quantities are derived for policy and food security assessment.

\subsection{Data}
\label{subsec:data}

The analysis uses two main data layers. \textbf{Micro data} come from a household food consumption survey underlying \texttt{data.csv}. The micro data record at-home and total (at-home plus out-of-home) consumption by food category, household income, and standard covariates. In the current replication package, the micro table contains 47{,}426 household--wave observations from 6{,}825 households, covering 11 provinces over the survey years 2004--2011 (variables \texttt{hhid}, \texttt{T1}, \texttt{wave}). Because the underlying microdata are access-restricted, we do not redistribute them; nonetheless, we document variable definitions and coverage in this paper and provide all downstream processed outputs. \textbf{Macro data} are province--year aggregates (e.g., \texttt{data2012.csv}), constructed from official provincial and national statistics; they provide average income and socio-economic covariates (e.g., urbanization rate, aging rate, Engel coefficient, food price index) by province and year but do not contain consumption variables. \textbf{Production data} (e.g., \texttt{data\_production1/2/3.csv}) are merged by province and year from the relevant statistical sources. \textbf{Per-capita at-home consumption} by province and category (e.g., \texttt{data\_q.csv}) and \textbf{population} (file name as in the replication package) for provincial and national weighting are obtained from the same or complementary statistical sources. Our target reporting period (2015--2024) is covered by the macro tables and output CSVs; prediction for those years is therefore an extrapolation from micro-based relationships combined with macro covariates and the bridging steps described below.

\subsection{Data Difficulties}
\label{subsec:data-difficulties}

\subsubsection{Micro versus macro data structure}
\label{subsec:micro-macro}

We face a fundamental \textbf{data-level mismatch}. \textbf{Micro data} (household survey, e.g., \texttt{data.csv}) are at the individual or household level. They contain \emph{at-home} and \emph{total} (at-home plus out-of-home) consumption by category, so that the at-home consumption ratio (at-home / total) and thus the out-of-home coefficient (total / at-home) can be computed for each observation. They also include \textbf{individual income} (\texttt{indinc}) and covariates such as province (\texttt{T1}), year (\texttt{wave}), and socio-economic variables. In contrast, \textbf{macro data} (province--year aggregates, e.g., \texttt{data2012.csv}) provide \textbf{average income} (\texttt{income}) and macro covariates (e.g., urbanization rate, aging rate, Engel coefficient, food price index, and other regional variables), but \textbf{no consumption variables} and \textbf{no individual income}. Out-of-home consumption is not directly observed at the macro level; it must be inferred by bridging micro and macro data.

\subsubsection{Income variable mismatch}
\label{subsec:income-mismatch}

Models trained on micro data use \textbf{individual income} (\texttt{indinc}) as a predictor of the at-home consumption ratio. At prediction time we only have \textbf{province--year average income} (\texttt{income}) for the macro table. The macro table has no income \emph{distribution} and no \texttt{indinc}. Applying the trained model directly to macro data would therefore require an \texttt{indinc} column that does not exist. Simply substituting \texttt{income} for \texttt{indinc} would ignore within-province income inequality and the non-linear relationship between income and consumption structure, and would be inconsistent with the training feature set. Thus, a method is needed to \textbf{match the income distribution} from the micro data to the macro prediction context and to \textbf{generate synthetic individual income values} for macro observations so that the same feature space is used in training and prediction.

\subsubsection{Missing provinces and years}
\label{subsec:missing-provinces}

Micro survey data do not cover all provinces or all years uniformly. For some province--year cells there are \textbf{no micro observations}; for others, consumption or ratios may be \textbf{missing} for certain categories (e.g., zero or unreported consumption). Macro prediction is required for all provinces and for years from 2015 onward. Therefore, (i) \textbf{missing at-home/total consumption or ratios} in the micro data must be handled (e.g., by imputation or by restricting to complete cases for robustness checks), and (ii) \textbf{provinces with no or insufficient micro data} require a way to obtain out-of-home (or at-home) coefficients---typically via \textbf{spatial interpolation} from provinces where predictions are available, possibly combined with auxiliary information such as provincial production shares.

\subsubsection{Other data limitations}
\label{subsec:other-data}

Additional constraints include: the need to align \textbf{production data} (e.g., \texttt{data\_production1/2/3.csv}) and \textbf{per-capita at-home consumption} (e.g., \texttt{data\_q.csv}) with province and year for use as features or in post-processing; the use of \textbf{population} (e.g., \texttt{procince\_pop.csv}) only at the national level for weighting; and standard data-quality issues (zero or negative consumption, ratio bounds, outliers). The research design must explicitly account for these so that estimates remain consistent and interpretable.

\subsection{Methodological Difficulties}
\label{subsec:methodological-difficulties}

\subsubsection{Why standard approaches are insufficient}
\label{subsec:why-standard-insufficient}

A \textbf{direct macro regression} of out-of-home consumption on province--year covariates is not feasible: macro data do not contain any measure of out-of-home (or total) consumption, so there is no outcome to regress. A \textbf{micro-only} analysis can describe ratios and coefficients for the sample but cannot produce province--year or national estimates for out-of-home consumption without a formal way to \textbf{scale from micro to macro} and to fill missing provinces and years. Simple \textbf{mean substitution} of \texttt{income} for \texttt{indinc} in prediction ignores the role of income distribution and is not consistent with the training specification. Ignoring missing provinces would leave gaps in the coverage required for policy; ignoring missing ratios would discard information and may bias estimates if missingness is not random. Thus, the methodology must address: (i) \textbf{distribution matching} so that macro prediction uses a synthetic income variable aligned with the micro distribution; (ii) \textbf{imputation or explicit handling of missing ratios} in the training data; (iii) \textbf{spatial interpolation} for provinces without micro-based predictions; (iv) \textbf{multi-category and grain-structure prediction} (four grain types plus other categories); and (v) \textbf{quantification of uncertainty} (e.g., bootstrap) rather than point estimates alone.

\subsubsection{Need for distribution matching, spatial interpolation, and ML prediction}
\label{subsec:need-for-methods}

\textbf{Distribution matching} is necessary because the model is trained on \texttt{indinc} and must be applied to macro rows that only have \texttt{income}. We therefore generate a synthetic micro-style income variable for macro rows using a rank-preserving quantile-scaling procedure (Section~\ref{sec:copula}): the shape of the micro income distribution is preserved while the mean is scaled to the macro \texttt{income}. \textbf{Spatial interpolation} (inverse-distance weighting, IDW) is needed to assign coefficients to provinces that lack micro data, using predictions from nearby provinces and, for grain structure, auxiliary production information. \textbf{Machine learning (ML) prediction} is adopted because the relationship between covariates (income, macro variables, production, province, year) and the at-home consumption ratio is likely nonlinear and may vary by category; a suite of ML models allows comparison and selection (including for imputation) and robustness checks. Together, these choices are dictated by the data: without distribution matching we cannot apply the micro-trained model to macro; without interpolation we cannot cover all provinces; without ML we may underfit complex, category-specific relationships.

\subsection{Research Strategy}
\label{subsec:research-strategy}

The strategy implemented in this project is structured in five main steps, aligned with the project documentation.

\subsubsection{Data preparation and income-distribution matching}
\label{subsec:copula}

Micro data (\texttt{data.csv}) and macro data (\texttt{data2012.csv}) are loaded; production and other auxiliary datasets are merged by province (\texttt{T1}) and year (\texttt{wave}). For macro prediction (from 2015 onward), we apply \textbf{income-distribution matching} once in the data preparation phase. The implementation (\texttt{match\_income\_copula} in \texttt{data\_preparation\_advanced.py}) uses a rank-preserving quantile-scaling procedure: it preserves the shape of the micro income distribution while scaling its mean to the macro \texttt{income}, generating \textbf{synthetic \texttt{indinc}} values for the macro table. Thus, the same feature set (including \texttt{indinc}) is available for both training (micro) and prediction (macro). This step is category-independent and is best interpreted as distribution (quantile) matching rather than a fully parametric Copula model.

\subsubsection{Imputation for missing ratios and model selection}
\label{subsec:imputation}

Where at-home or total consumption (and hence the at-home ratio) is missing in the micro data, we first evaluate \textbf{multiple ML models} (e.g., 12 models in the imputation selection step) for their imputation accuracy via cross-validation (e.g., MAE, RMSE, R²). The \textbf{best-performing model} is used to produce a single imputed ratio dataset (\texttt{data/imputed\_ratios\_best.csv}), which is then used consistently in the main prediction pipeline. This avoids ad hoc choice of imputation and allows a \textbf{robustness check} by comparing results with and without imputation (main pipeline with imputation vs. robust pipeline without imputation).

\subsubsection{Model training and prediction by category}
\label{subsec:training-prediction}

Models are \textbf{trained on micro data}: features include \texttt{indinc} (and other constructed features such as macro and production variables aligned to the micro level), and the target is the \textbf{at-home consumption ratio} (at-home / total) for each of the 15 categories. The \textbf{out-of-home coefficient} is then defined as the reciprocal of the predicted ratio. For \textbf{grain structure}, the four grain types (rice, wheat, beans, coarse grains) are modeled separately (e.g., shares derived from micro data), and missing provinces are handled by spatial interpolation combined with production shares. For \textbf{out-of-home coefficients}, the same 12 ML models (LightGBM, XGBoost, CatBoost, Random Forest, Lasso, Ridge, MLP, TabNet, TabM, TabPFN, ResNet, FT-Transformer) are run; where a province has no micro-based prediction, \textbf{spatial interpolation} (inverse-distance weighting, IDW) is applied using predicted coefficients from other provinces and optionally production weights.

\subsubsection{Post-processing and aggregation}
\label{subsec:postprocess}

Using predicted out-of-home coefficients, grain structure, per-capita at-home consumption (\texttt{data\_q.csv}), and population (\texttt{procince\_pop.csv}), \textbf{provincial} per-capita at-home and total (at-home plus out-of-home) consumption are computed by category; for grains, provincial at-home consumption is allocated by grain-type shares before applying the coefficient. \textbf{National} aggregates are obtained by population-weighted averages (e.g., national out-of-home coefficient and per-capita consumption). Outputs include \texttt{results\_province\_<model>.csv} and \texttt{results\_national\_<model>.csv}, and a \texttt{model\_comparison\_national.csv} for cross-model comparison.

\subsubsection{Robustness and uncertainty}
\label{subsec:robustness}

\textbf{Robustness} is assessed by running the same models \textbf{without imputation} (only observations with non-missing ratios), producing \texttt{predictions\_<model>\_robust.csv} and corresponding provincial/national results, and comparing them to the main (imputation) pipeline. \textbf{Uncertainty} is quantified via \textbf{Bootstrap}: multiple runs with different random seeds yield a distribution of predictions, from which means and percentile intervals (e.g., 2.5\%--97.5\%) are computed (\texttt{predictions\_<model>\_bootstrap.csv}).

\subsection{Justification of the Approach}
\label{subsec:justification}

Given the existing data, this strategy is \textbf{scientifically justified} and \textbf{not easily replaceable} for the following reasons.

\begin{itemize}
  \item \textbf{No alternative to bridging micro and macro.} Macro data do not contain out-of-home (or total) consumption; the only source of consumption structure is the micro survey. Therefore, any province--year or national estimate must use the micro data to learn the relationship between covariates and the at-home ratio, then apply it to macro cells. The chosen pipeline does exactly that: train on micro, predict on macro after aligning the feature space.

  \item \textbf{Income-distribution matching is required by the feature set.} The model is trained with \texttt{indinc}; macro rows have only \texttt{income}. Without generating a consistent \texttt{indinc} for macro (via distribution matching), prediction is either undefined or relies on an incorrect substitution. The rank-preserving scaling procedure preserves distributional information from the micro data and scales to the macro mean, making the application of the micro-trained model to macro data coherent.

  \item \textbf{Imputation and spatial interpolation are necessitated by coverage.} Missing ratios in the micro data and missing provinces in the prediction target cannot be ignored if the goal is full provincial and national coverage. Imputation (with a selected model and a no-imputation robustness check) and spatial interpolation (with production information where relevant) are therefore not optional but part of a minimal feasible design.

  \item \textbf{Multiple ML models and checks support transparency.} Using 12 ML models allows comparison of performance and stability across categories; robustness (with/without imputation) and Bootstrap intervals make the dependence on imputation and the uncertainty of estimates explicit. This is appropriate for policy-oriented work targeting top agricultural economics journals.

\end{itemize}

In sum, the research design is determined by \textbf{data availability} and the \textbf{research question}: the combination of micro-based training, income-distribution matching, imputation, spatial interpolation, multi-model prediction, and robustness and uncertainty checks is the scientifically appropriate and effectively unique way to estimate China's out-of-home food consumption at national and provincial levels by category with the existing project data and documentation.


% ========== 4. Methods (PhD3) ==========
% Research Methods Section
% PhD Student 3 — Machine learning models, pipeline, and applied tests
% Paper: China's out-of-home food consumption
% No \begin{document}; for inclusion in main manuscript.

\section{Research Methods}
\label{sec:methods}

\subsection{Prediction Target and Outcome Variable}

The training and prediction target is the \emph{at-home consumption ratio} (ratio of at-home to total at-home plus out-of-home consumption), defined as
\begin{equation}
  \text{ratio} = \frac{\text{at-home}}{\text{at-home} + \text{out-of-home}} \in (0,\,1].
\end{equation}
The \emph{out-of-home consumption coefficient} is then defined as the reciprocal of this ratio:
\begin{equation}
  \text{out-of-home coefficient} = \frac{1}{\text{ratio}} = \frac{\text{at-home} + \text{out-of-home}}{\text{at-home}}.
\end{equation}
Total (at-home plus out-of-home) consumption is obtained in post-processing as
\begin{equation}
  \text{at-home} + \text{out-of-home} = \text{at-home} \times \text{out-of-home coefficient}.
\end{equation}
This relation is applied in the post-processing step (\texttt{postprocess\_predictions.py}) to derive provincial and national consumption quantities from predicted out-of-home coefficients and at-home consumption from \texttt{data\_q.csv}.

For interpretation, the coefficient is a \emph{scaling factor} that maps at-home quantities into total quantities under aligned definitions. The implied out-of-home share of total consumption is
\begin{equation}
  \text{share}_{\text{out}} = \frac{\text{out-of-home}}{\text{total}} = 1 - \frac{1}{\text{out-of-home coefficient}}.
\end{equation}
We emphasize that this paper estimates scaling factors for accounting and projection, not behavioral parameters or causal effects.

Predictions are produced for \textbf{15 food categories}: four grain types---paddy rice (\texttt{daogu}), wheat (\texttt{xiaomai}), beans (\texttt{doulei}), and coarse grains (\texttt{zaliang})---plus vegetables (\texttt{shucai}); meat (pork \texttt{zhurou}, beef \texttt{niurou}, mutton \texttt{yangrou}, poultry \texttt{qinlei}); aquatic products (\texttt{shuichanpin}); eggs (\texttt{danlei}); dairy (\texttt{nailei}); fruit (\texttt{guaguo}); edible oil (\texttt{youliao}); and sugar (\texttt{tang}).

\subsection{Features and Data Preparation}

Features are constructed in \texttt{data\_preparation\_advanced.py} in the function \texttt{load\_and\_prepare\_data\_advanced}. The feature vector used for both training and prediction is
\begin{equation}
  X = [\texttt{indinc},\, \texttt{income},\, \texttt{urbanrate},\, \texttt{oldrate},\, \texttt{engel},\, \texttt{foodpindex},\, \texttt{zyywsr},\, \texttt{cyyye},\, \texttt{cyyysr},\, \text{production variables},\, \texttt{T1},\, \texttt{wave}].
\end{equation}

\paragraph{Micro-level.}
\texttt{indinc} is household income aligned with the macro prediction table. When income-distribution matching is enabled, it is generated by \texttt{match\_income\_copula} so that the (rank-preserving) distribution of income in the macro scenario matches the micro survey distribution after scaling to macro means (see Section~\ref{sec:copula}).

\paragraph{Macro-level (province $\times$ year).}
Macro features include \texttt{income}, \texttt{urbanrate}, \texttt{oldrate}, \texttt{engel}, \texttt{foodpindex}, \texttt{zyywsr}, \texttt{cyyye}, \texttt{cyyysr}, and production variables. Production variables are merged from \texttt{data\_production1.csv}, \texttt{data\_production2.csv}, and \texttt{data\_production3.csv} by \texttt{wave} and \texttt{T1}; they include \texttt{qliangshi}, \texttt{qdaogu}, \texttt{qxiaomai}, \texttt{qyumi}, \texttt{qdoulei}, \texttt{qshulei}, \texttt{qyouliao}, \texttt{qshuiguo}, \texttt{qroulei}, \texttt{qzhurou}, \texttt{qniurou}, \texttt{qyangrou}, \texttt{qnailei}, \texttt{qqindan}, and \texttt{qshuichan}.

\paragraph{Identifiers.}
\texttt{T1} is the province code and \texttt{wave} is the year. Prediction is performed only for \texttt{wave} $\geq$ 2015 (\texttt{MIN\_PREDICT\_YEAR}).

Training uses the micro dataset \texttt{data.csv} with the above features and category-specific \texttt{ratio\_*} (at-home over at-home plus out-of-home) as the target, without truncation. Prediction uses the macro table \texttt{data2012.csv} and production tables to build the province$\times$year feature matrix for each category.

\subsection{Machine Learning Models}
\label{sec:models}

We use \textbf{12} machine learning models for coefficient prediction and, where applicable, for imputation. Each is implemented in a dedicated script (\texttt{predict\_<model>.py} or \texttt{predict\_fttransformer\_advanced.py}).

\paragraph{Tree-based and linear.}
\begin{itemize}
  \item \textbf{LightGBM}: gradient boosting (LGBMRegressor); fixed hyperparameters; used for coefficient prediction and in imputation selection.
  \item \textbf{XGBoost}: extreme gradient boosting; time-series cross-validation (TimeSeriesSplit, 3 splits) in training; coefficient prediction and imputation.
  \item \textbf{CatBoost}: gradient boosting with categorical support; TimeSeriesSplit (3 splits); coefficient prediction and imputation.
  \item \textbf{Random Forest}: ensemble of trees; coefficient prediction and imputation.
  \item \textbf{Lasso}: L1-regularized linear regression; coefficient prediction and imputation.
  \item \textbf{Linear (Ridge)}: L2-regularized linear regression; coefficient prediction and imputation.
\end{itemize}

\paragraph{Deep tabular.}
\begin{itemize}
  \item \textbf{MLP}: multilayer perceptron (\texttt{sklearn.neural\_network.MLPRegressor}); early stopping; fallback for TabNet/TabPFN when unavailable; coefficient prediction and imputation.
  \item \textbf{TabNet}: attention-based tabular model (pytorch-tabnet TabNetRegressor); \texttt{patience} for early stopping; fallback to MLP if unavailable; coefficient prediction and imputation.
  \item \textbf{TabM}: parameter-efficient MLP-style model (yandex-research/tabm); current implementation may fall back to MLP; coefficient prediction and imputation.
  \item \textbf{TabPFN}: prior-fitted network (TabPFNRegressor); local checkpoint when available; fallback to MLP; coefficient prediction and imputation.
  \item \textbf{ResNet (tabular)}: residual network for tabular data; coefficient prediction and imputation.
  \item \textbf{FT-Transformer}: feature-tokenizer Transformer (PyTorch); \texttt{DirectRatioOptimizer} with Optuna for coefficient prediction; dedicated \texttt{FTRegressionModel} for imputation (home/total, no sigmoid). FT-Transformer uses its own imputation when selected as best; it is still evaluated in imputation selection via \texttt{evaluate\_imputation\_cv\_fttransformer} and can be chosen as the single best model.
\end{itemize}

All 12 models participate in the imputation model selection (\texttt{run\_imputation\_selection.py}); the best model (e.g., FT-Transformer, as in \texttt{data/imputation\_selection\_results.csv}) is used to produce the unified imputed sample \texttt{data/imputed\_ratios\_best.csv}.

\subsection{Pipeline Overview}
\label{sec:pipeline}

The workflow consists of the following steps.

\begin{enumerate}
  \item \textbf{Data preparation.}
  Load micro \texttt{data.csv}, macro \texttt{data2012.csv}, and production tables; merge production by \texttt{wave} and \texttt{T1}; restrict prediction to years $\geq$ 2015. Optionally apply income-distribution matching (rank-preserving scaling) to obtain \texttt{indinc} for the macro table. Output: prepared micro data, prediction feature matrix \texttt{X\_pred\_feats}, feature list \texttt{feature\_cols}, category map, and set of known provinces.

  \item \textbf{Imputation model selection} (\texttt{run\_imputation\_selection.py}).
  For each category, the 12 models (including FT-Transformer) are evaluated on \emph{observed} at-home/total data using 5-fold cross-validation. Each model predicts at-home and total consumption; the ratio is computed and compared to the observed ratio. Metrics (MAE, RMSE, R²) are averaged across categories per model; the model with the lowest average MAE is selected. The selected model is then used to impute missing at-home/total (and thus ratio) for all categories; results are saved as \texttt{data/imputed\_ratios\_best.csv} and the model name in \texttt{data/best\_imputation\_model.txt}. Aggregated metrics and ranks are written to \texttt{data/imputation\_selection\_results.csv}.

  \item \textbf{Coefficient prediction (main pipeline).}
  If \texttt{data/imputed\_ratios\_best.csv} exists, each model uses the unified imputed sample (\texttt{ratio\_filled\_*}) to train the coefficient model and to predict the at-home consumption ratio for the macro grid. Otherwise, each model imputes internally. For each category, the out-of-home coefficient is obtained as the reciprocal of the predicted ratio. Missing provinces are filled via spatial interpolation (inverse-distance weighting, IDW). Outputs: \texttt{predictions\_<model>.csv} (out-of-home coefficients by province and year).

  \item \textbf{Grain structure.}
  Four grain shares (paddy, wheat, beans, coarse grains) are predicted from micro data (e.g., paddy from total rice/0.7, then shares of total grain mass). For provinces missing in the micro data, shares are obtained by spatial interpolation (IDW) combined with provincial production shares. Results are saved as \texttt{grain\_structure\_predictions\_<model>.csv} and used in post-processing.

  \item \textbf{Post-processing} (\texttt{postprocess\_predictions.py}).
  Using predicted out-of-home coefficients, grain structure, \texttt{data\_q.csv} (provincial per-capita at-home consumption), and \texttt{procince\_pop.csv} (population): provincial per-capita at-home consumption for each grain type is \texttt{q\_liangshi} times the corresponding share; provincial per-capita at-home plus out-of-home is per-capita at-home times the out-of-home coefficient. For non-grain categories, provincial per-capita at-home is taken from \texttt{data\_q.csv} and multiplied by the out-of-home coefficient to get per-capita at-home plus out-of-home. National out-of-home coefficients and per-capita consumption are computed as population-weighted averages across provinces. Outputs: \texttt{results\_province\_<model>.csv}, \texttt{results\_national\_<model>.csv}.

  \item \textbf{Robustness (no imputation).}
  Each model is also run without imputation: only samples with observed ratio (and consumption $>0$) are used for training and prediction. Results are saved as \texttt{predictions\_<model>\_robust.csv} and can be compared to the main pipeline.

  \item \textbf{Bootstrap prediction intervals} (\texttt{run\_bootstrap.py}).
  For a given model, prediction is run multiple times (default 30) with different random seeds (\texttt{PREDICT\_SEED}). For each (Province, Year) and each category, the mean and 2.5\% and 97.5\% percentiles are computed and saved as \texttt{predictions\_<model>\_bootstrap.csv} (columns \texttt{*\_Mean}, \texttt{*\_Lower}, \texttt{*\_Upper}).
\end{enumerate}

\subsection{Supporting Methods}

\subsubsection{Income-distribution matching (rank-preserving scaling)}
\label{sec:copula}

Micro data contain household income (\texttt{indinc}), while the macro prediction table has only province--year mean income (\texttt{income}). To align the training and prediction feature spaces, we apply income-distribution matching implemented in \texttt{match\_income\_copula} in \texttt{data\_preparation\_advanced.py}. Despite the function name, the procedure is best interpreted as a \emph{rank-preserving quantile-scaling} approach rather than a fully parametric Copula model: for each province--year in the macro table, the micro income distribution is taken from the empirical distribution of household income in the corresponding micro sample where available (or from the pooled micro sample otherwise). A scale factor is computed so that the mean of the scaled distribution equals the macro mean income for that province--year (scale factor = macro mean / micro mean); empirical quantiles of the micro distribution are then multiplied by this factor to produce synthetic \texttt{indinc} values for the macro table. This step runs once during data preparation and is independent of food category.

\subsubsection{Spatial interpolation for missing provinces}

A province is classified as \emph{missing} if it has no micro data in the survey (i.e., it does not appear in the micro sample at all), as opposed to having no data only for a particular year. For each year, coefficients for missing provinces are filled using spatial interpolation from provinces that do have micro data. \textbf{Out-of-home coefficients:} \texttt{apply\_kriging\_interpolation} in \texttt{data\_preparation\_advanced.py} uses inverse-distance weighting (IDW; distance-based weights) to fill missing provinces for the out-of-home coefficient, by year. \textbf{Grain structure:} \texttt{apply\_grain\_structure\_interpolation} uses IDW (inverse-distance squared weights) combined with the province’s production share of total grain for that type; the final share is a blend of the spatial interpolant and the production-based share. An alternative IDW-only path is available in \texttt{apply\_spatial\_interpolation} for scripts that call it.

\subsubsection{Grain structure (four shares)}

Grain is split into four types: paddy (稻谷), wheat (小麦), beans (豆类), and coarse grains (杂粮). From micro data, paddy mass is derived as total rice/0.7; total grain mass is the sum of the four components. Each share is modeled separately; for provinces without micro data, the share is obtained by IDW plus production share as above. The four shares are normalized to sum to one and written to \texttt{grain\_structure\_predictions\_<model>.csv}. In post-processing, provincial per-capita at-home consumption for each grain type is \texttt{q\_liangshi} times the corresponding share; the out-of-home coefficient for that type is then applied to obtain per-capita at-home plus out-of-home consumption.

\subsection{Evaluation and Regularization}

\begin{itemize}
  \item \textbf{Imputation selection.} 5-fold cross-validation on observed (at-home, total) pairs; metrics MAE, RMSE, R²; selection by lowest average MAE across categories; 12 models compared; results in \texttt{data/imputation\_selection\_results.csv} and \texttt{data/imputation\_evaluation\_detail.csv}.
  \item \textbf{Robustness.} Training and prediction using only observed ratio (no imputation), yielding \texttt{predictions\_<model>\_robust.csv} for comparison with the main pipeline.
  \item \textbf{Bootstrap.} Multiple runs with different seeds; mean and 2.5\%/97.5\% percentiles per (Province, Year) and category in \texttt{predictions\_<model>\_bootstrap.csv}.
  \item \textbf{Cross-validation.} 5-fold CV in imputation selection; TimeSeriesSplit (3 splits) in coefficient training for several models (e.g., XGBoost, CatBoost, FT-Transformer).
  \item \textbf{Early stopping.} Used where applicable: e.g., MLP (\texttt{early\_stopping=True}), TabNet (\texttt{patience}), FT-Transformer (validation loss and \texttt{patience}); reduces overfitting.
  \item \textbf{Regularization.} L1/L2 in Lasso/Ridge; fixed or Optuna-tuned hyperparameters for tree and deep models (FT-Transformer uses Optuna for coefficient prediction).
\end{itemize}

All methodological claims above are aligned with the implementation in \texttt{data\_preparation\_advanced.py}, \texttt{run\_imputation\_selection.py}, \texttt{run\_bootstrap.py}, \texttt{postprocess\_predictions.py}, and the \texttt{predict\_*.py} scripts, and with the project documentation.


% ========== 5. Results (PhD4) ==========
% Results section -- PhD4 (Results and Discussion specialist)
% All numbers traceable to CSVs in project path /usr/fafh.
% Source: model_comparison_national.csv, results_national_*.csv, data/imputation_selection_results.csv,
%         results_national_*_robust.csv, model_accuracy_summary.csv.

\section{Results}
\label{sec:results}

\subsection{Main Estimates of National Out-of-Home Consumption Coefficients}

We estimate national out-of-home consumption coefficients for 15 food categories and 10 years (2015--2024) using 12 prediction models. The coefficient for each category is the ratio of total (at-home plus out-of-home) consumption to at-home consumption; a value above one indicates that consumption away from home raises total consumption above the at-home level. All estimates in this subsection are taken from \texttt{model\_comparison\_national.csv} and the \texttt{results\_national\_*.csv} files (e.g., \texttt{results\_national\_lightgbm.csv}, \texttt{results\_national\_tabpfn.csv}, \texttt{results\_national\_fttransformer.csv}).

Across models and years, grain categories (rice \textit{daogu}, wheat \textit{xiaomai}, legumes \textit{doulei}, coarse grains \textit{zaliang}) show coefficients typically between about 1.05 and 1.35. For example, in 2020 LightGBM yields 1.128 for rice, 1.153 for wheat, 1.182 for legumes, and 1.153 for coarse grains (\texttt{model\_comparison\_national.csv}, rows for Year=2020, model=lightgbm). Vegetables (\textit{shucai}) are generally in the 1.05--1.19 range (e.g., LightGBM 2020: 1.069; TabPFN 2024: 1.187). Meat and animal products show larger multipliers: pork (\textit{zhurou}) often between 1.16 and 1.38, poultry (\textit{qinlei}) and beef (\textit{niurou}) in similar or higher ranges depending on model and year. Egg (\textit{danlei}), dairy (\textit{nailei}), fruits (\textit{guaguo}), cooking oil (\textit{youliao}), and sugar (\textit{tang}) coefficients are usually closer to one or modestly above (e.g., oil and sugar near 1.00--1.02 in many specifications).

The tree-based and neural/table models give broadly consistent orders of magnitude, but levels differ. LightGBM tends to produce somewhat lower coefficients for grains and meat in later years (e.g., 2024 rice 1.110, pork 1.168), whereas TabPFN and FT-Transformer often give higher values (e.g., TabPFN 2024 rice 1.233, pork 1.376; FT-Transformer 2024 rice 1.242, pork 1.228). Linear and Lasso specifications sit in the middle for many categories; some flexible models (e.g., XGBoost, ResNet, linear in a few categories) exhibit unstable or extreme point estimates in particular year--category cells and are therefore less reliable for policy interpretation. We emphasize LightGBM, TabPFN, and FT-Transformer for their combination of prediction accuracy and stability across categories: TabPFN attains the best prediction MAE, FT-Transformer is the best imputation model and ranks first on RMSE for the prediction task, and LightGBM provides stable point estimates consistent with the tabular-deep-learning models.

Table~\ref{tab:national_coef_summary} reports national out-of-home consumption coefficients for selected food categories and years from LightGBM and TabPFN. Every entry is taken from \texttt{model\_comparison\_national.csv} (columns \texttt{全国户外消费系数\_daogu}, \texttt{全国户外消费系数\_xiaomai}, \texttt{全国户外消费系数\_zhurou}, \texttt{全国户外消费系数\_shucai}).

\begin{table}[htbp]
\centering
\caption{National out-of-home consumption coefficients, selected years and food categories (LightGBM and TabPFN). Source: \texttt{model\_comparison\_national.csv}.}
\label{tab:national_coef_summary}
\begin{tabular}{llcccc}
\toprule
Model    & Year & Rice (daogu) & Wheat (xiaomai) & Pork (zhurou) & Vegetables (shucai) \\
\midrule
LightGBM & 2015 & 1.135 & 1.177 & 1.234 & 1.073 \\
         & 2020 & 1.128 & 1.153 & 1.181 & 1.069 \\
         & 2024 & 1.110 & 1.175 & 1.168 & 1.062 \\
\midrule
TabPFN   & 2015 & 1.139 & 1.213 & 1.313 & 1.124 \\
         & 2020 & 1.157 & 1.161 & 1.327 & 1.129 \\
         & 2024 & 1.233 & 1.149 & 1.376 & 1.187 \\
\bottomrule
\end{tabular}
\end{table}

\subsection{Imputation Model Selection}

Missing values in the out-of-home consumption ratio data are imputed using a supervised learning procedure. Model selection is based on hold-out performance on the imputation task. Metrics and ranks are reported in \texttt{data/imputation\_selection\_results.csv}.

The best-performing imputation model is FT-Transformer (\texttt{is\_best=True}): MAE = 0.1704, RMSE = 0.2978, with rank 1. ResNet Tabular ranks second (MAE 0.1793, RMSE 0.3054), followed by MLP, TabM, TabPFN, Random Forest, CatBoost, LightGBM, Linear, Lasso, XGBoost, and TabNet (worst MAE 0.2512). The negative $R^2$ values in that file reflect out-of-sample evaluation and indicate that the imputation task is difficult; MAE and RMSE are the primary criteria for selecting the FT-Transformer as the best imputation model. All downstream national and provincial results that use imputation therefore rely on the FT-Transformer-imputed ratios unless otherwise stated.

\subsection{Robustness: Main Versus No-Imputation (Robust) Specifications}

To assess sensitivity to imputation, we re-estimate national coefficients using the same prediction models but without imputing missing ratios (robust runs). The main results are in \texttt{results\_national\_<model>.csv}; the no-imputation counterparts are in \texttt{results\_national\_<model>\_robust.csv}.

For LightGBM, the main and robust coefficients are close in level and trend. In 2015, the main rice coefficient is 1.135 and the robust 1.141 (\texttt{results\_national\_lightgbm.csv} and \texttt{results\_national\_lightgbm\_robust.csv}, column \texttt{全国户外消费系数\_daogu}); in 2024, main is 1.110 and robust 1.140. Pork in 2024 is 1.168 (main) vs.\ 1.129 (robust). Omitting imputation thus slightly raises the rice coefficient and slightly lowers the pork coefficient for this model.

For TabPFN, the robust coefficients are generally somewhat higher for grains and some other categories. In 2015, rice is 1.139 (main) vs.\ 1.221 (robust); in 2024, 1.233 (main) vs.\ 1.255 (robust). Vegetables in 2024 are 1.187 (main) vs.\ 1.173 (robust). For FT-Transformer, 2015 rice is 1.206 (main) vs.\ 1.245 (robust) and 2024 rice 1.242 (main) vs.\ 1.263 (robust). The direction and size of the main--robust gap vary by model and category, but in no case do we find a qualitative reversal: out-of-home consumption continues to add to total consumption (coefficients remain above one), and the magnitudes remain in a comparable range. We conclude that the main findings are robust to excluding imputation, with the caveat that point estimates can shift by a few percentage points depending on the model and food category.

\subsection{Model Performance on the Prediction Task}

Prediction performance for the at-home consumption ratio (the training target) is summarized in \texttt{model\_accuracy\_summary.csv}. These metrics reflect internal predictive accuracy on the microdata evaluation design used in the project pipeline; they do not constitute external validation against independent macro-level measurements, which are unavailable for China at the required resolution. In the summary file, TabPFN attains the lowest MAE (0.0797), followed by TabM (0.0852), FT-Transformer (0.0865), and LightGBM (0.0872). RMSE rankings are similar, with FT-Transformer having the lowest RMSE (0.1432). Category-level metrics in \texttt{model\_accuracy\_detail.csv} show that performance varies by food category; for instance, FT-Transformer achieves the highest $R^2$ for rice (q\_daogu) among the models considered. Taken together with cross-model stability in the national and provincial outputs, these results motivate our emphasis on LightGBM, TabPFN, and FT-Transformer when presenting headline estimates.

\subsection{Uncertainty Quantification}

Bootstrap prediction intervals are produced by the project pipeline (\texttt{run\_bootstrap.py}): the prediction step is run repeatedly with different random seeds, and the 2.5\%--97.5\% percentiles across runs yield 95\% intervals for each province--year--category cell. Full bootstrap-based 95\% intervals for all models and categories are stored in the replication outputs \texttt{predictions\_<model>\_bootstrap.csv} (columns \texttt{*\_Mean}, \texttt{*\_Lower}, \texttt{*\_Upper}). These intervals confirm that the main qualitative conclusion---coefficients above one for all categories---is unchanged when uncertainty is accounted for, while also making transparent the residual model and sampling uncertainty for policy applications.

\subsection{Provincial Heterogeneity}

Provincial out-of-home consumption coefficients are estimated for all 12 models and stored in \texttt{results\_province\_*.csv}. Table~\ref{tab:provincial_heterogeneity} summarizes coefficient heterogeneity across selected provinces and years for the TabPFN model (source: \texttt{results\_province\_tabpfn.csv}, columns \texttt{q\_daogu\_coef}, \texttt{q\_zhurou\_coef}, \texttt{q\_shucai\_coef}). Province codes follow the project convention (GB/T 2260): 11 = Beijing, 13 = Hebei, 21 = Liaoning. Coefficients vary substantially by province: for example, in 2024 rice (daogu) ranges from about 1.13 (province 21) to 1.53 (province 11), and pork (zhurou) from about 1.30 to 1.38. Full provincial results for all models, years, and categories are available in the replication outputs \texttt{results\_province\_*.csv}.

\begin{table}[htbp]
\centering
\caption{Provincial out-of-home consumption coefficients (TabPFN), selected provinces, years and categories. Source: \texttt{results\_province\_tabpfn.csv}.}
\label{tab:provincial_heterogeneity}
\begin{tabular}{llcccc}
\toprule
Province & Year & Rice (daogu) & Wheat (xiaomai) & Pork (zhurou) & Vegetables (shucai) \\
\midrule
11 & 2020 & 1.428 & 1.160 & 1.279 & 1.272 \\
11 & 2024 & 1.533 & 1.208 & 1.377 & 1.397 \\
21 & 2020 & 1.090 & 1.049 & 1.246 & 1.110 \\
21 & 2024 & 1.129 & 1.051 & 1.300 & 1.149 \\
13 & 2020 & 1.197 & 1.142 & 1.326 & 1.149 \\
13 & 2024 & 1.281 & 1.143 & 1.375 & 1.216 \\
\bottomrule
\end{tabular}
\end{table}


% ========== 6. Discussion (PhD4) ==========
% Discussion section -- PhD4 (Results and Discussion specialist)
% All quantitative claims traceable to CSVs in /usr/fafh.
% Sources: results_national_*.csv, model_comparison_national.csv, predictions_*.csv, grain_structure_predictions_*.csv.

\section{Discussion}
\label{sec:discussion}

\subsection{Main Out-of-Home Consumption Levels and Coefficients}

Our estimates imply that food away from home raises total consumption above at-home levels for all major food categories, with coefficients (ratio of total to at-home consumption) consistently above one. The magnitudes depend on the model and the food category. Using the TabPFN national results (\texttt{results\_national\_tabpfn.csv}) as a central reference, Table~\ref{tab:tabpfn_2024_national} reports the national out-of-home consumption coefficients for all 15 food categories in 2024. The corresponding at-home and total (at-home plus out-of-home) consumption levels (in kg per capita) are given in the same file (e.g., columns \texttt{户内消费量\_*}, \texttt{户内外消费量\_*}); for example, rice at-home consumption in 2024 is 54.24~kg per capita and total (at-home plus out-of-home) is 66.57~kg per capita. LightGBM (\texttt{results\_national\_lightgbm.csv}) yields somewhat lower coefficients for grains and meat in later years (e.g., 2024 rice 1.110, pork 1.168), while the best imputation model FT-Transformer (\texttt{results\_national\_fttransformer.csv}) gives values between LightGBM and TabPFN (e.g., 2024 rice 1.242, pork 1.228). The pattern is consistent: out-of-home consumption adds non-negligibly to total demand, with the largest relative contributions for poultry, pork, coarse grains, and legumes, and smaller relative contributions for staples such as cooking oil and sugar.

\begin{table}[htbp]
\centering
\caption{National out-of-home consumption coefficients, TabPFN, 2024. Source: \texttt{results\_national\_tabpfn.csv}, row Year=2024.}
\label{tab:tabpfn_2024_national}
\begin{tabular}{llc}
\toprule
Category (pinyin) & English & Coefficient \\
\midrule
daogu & Rice & 1.233 \\
xiaomai & Wheat & 1.149 \\
doulei & Legumes & 1.257 \\
zaliang & Coarse grains & 1.417 \\
shucai & Vegetables & 1.187 \\
zhurou & Pork & 1.376 \\
niurou & Beef & 1.369 \\
yangrou & Mutton & 1.077 \\
qinlei & Poultry & 1.753 \\
shuichanpin & Aquatic products & 1.458 \\
danlei & Eggs & 1.298 \\
nailei & Dairy & 1.160 \\
guaguo & Fruits & 1.123 \\
youliao & Cooking oil & 1.016 \\
tang & Sugar & 1.173 \\
\bottomrule
\end{tabular}
\end{table}

\subsection{Comparison with FAO, OECD, IFPRI and Related Literature}

International organizations and the literature on China's food consumption rarely report out-of-home consumption coefficients in the same form as ours (ratio of total to at-home consumption by food category). Definitions, coverage, and years differ, so direct numerical comparison is limited; we highlight conceptual links and where our estimates can inform or complement existing work.

\textbf{How to use the coefficients with external totals.} When an external source reports total consumption (or food-use) aggregates without an at-home/out-of-home split, our coefficient provides an accounting identity to recover implied components under aligned definitions. For a given product and year, if an external statistic provides a total per-capita quantity \(Q^{\text{total}}\), then implied at-home consumption is \(Q^{\text{home}} = Q^{\text{total}}/\text{coef}\), and the implied out-of-home share is \(1 - 1/\text{coef}\). This comparison should be undertaken only when product definitions and boundaries are aligned (e.g., edible-weight consumption versus supply-use aggregates), and we therefore emphasize conceptual complementarity rather than direct numeric equality across sources.

\textbf{FAO.} FAO Food Balance Sheets (FBS) and statistics focus on total food availability or supply (e.g., food supply quantity, dietary energy supply) at the national level and do not separate at-home and out-of-home consumption. FAO data are therefore not directly comparable to our coefficients, which are defined as total-to-at-home ratios. Our estimates can nevertheless help interpret total demand: if a share of consumption is eaten away from home and our coefficients are above one, then demand for primary products (e.g., grain for direct consumption and for feed and processing) is higher than what at-home consumption alone would imply. FAO FBS figures on per capita food use in China can be combined with our coefficients to back out implied at-home vs.\ out-of-home shares, subject to alignment of product definitions and years.

\textbf{Boundary alignment and remaining mismatches.} Even when using the accounting identity above, careful boundary alignment is required. Key sources of mismatch include: (i) supply--use aggregates versus edible-weight consumption quantities; (ii) inclusion of feed, processing, losses, and stock changes in supply-use accounts; (iii) treatment of institutional catering and business-provided meals; and (iv) differences in product aggregation (e.g., cereals versus rice). We therefore interpret comparisons to FAO/OECD/IFPRI primarily as conceptual cross-checks and as a demonstration of how our coefficients can be used to construct consistent at-home versus total splits when definitions are aligned.

\textbf{OECD.} The OECD tracks food consumption and agricultural outlooks; some OECD work discusses eating out and convenience foods in China but typically at an aggregate or value share level rather than physical quantities by food category. Our category-level coefficients (e.g., for rice, wheat, meat, vegetables) offer a quantity-based complement: they show by how much total consumption of each category exceeds at-home consumption, which can support OECD-style demand and trade scenario analysis when definitions and time coverage are aligned.

\textbf{IFPRI.} IFPRI has published on dietary change, food security, and consumption patterns in China, including shifts toward more animal products and processed foods. Again, few IFPRI reports provide out-of-home consumption coefficients in our exact form. Our results are consistent with a narrative of rising food away from home: coefficients above one across categories and the relative size of coefficients for meat and poultry align with the importance of dining out in urban and increasingly rural diets. Discrepancies with any IFPRI (or other) estimates would be expected where definitions differ (e.g., coverage of institutional catering, definition of “at home”) or where years and data sources differ.

\textbf{Other literature.} Studies on China's food away from home often use household survey data and focus on expenditure shares or frequency of eating out rather than physical consumption coefficients by commodity. Our machine-learning-based coefficients, derived from survey and auxiliary data and applied to national (and provincial) totals, are a different object: they yield total-to-at-home ratios by food category and year. As such, they are best compared to studies that estimate similar ratios or that decompose total consumption into at-home and out-of-home components. Where such studies exist, we would expect broad agreement in direction (out-of-home adding to total demand) and possible differences in level due to sample, period, and method. Our robustness checks (main vs. no-imputation specifications) suggest that our central conclusions are not driven by the choice of imputation.

\subsection{Implications for China's Food Security}

Our estimates have direct implications for the assessment of China's grain and food demand and hence for food security and self-sufficiency.

\textbf{Total grain demand.} Because the out-of-home consumption coefficients for grains (rice, wheat, legumes, coarse grains) are above one, total grain consumption for direct human use is higher than at-home consumption alone. The exact increment depends on the model: for 2024, the rice coefficient ranges from about 1.11 (LightGBM) to 1.23 (TabPFN) to 1.24 (FT-Transformer) in \texttt{model\_comparison\_national.csv}. Applying these coefficients to at-home consumption implies that a non-negligible share of grain demand is attributable to eating away from home. Any projection or policy analysis that ignores this component will understate total grain demand for human consumption.

\textbf{Demand structure.} The coefficients for meat (pork, beef, poultry), aquatic products, eggs, and vegetables are also above one and often larger than for some staples. This implies that out-of-home consumption is relatively important for protein-rich and higher-value foods. Demand for feed grains and for processing (e.g., flour for catering) is therefore linked to both at-home and out-of-home consumption; our coefficients help quantify the extra demand coming from the latter. Discrepancies between our estimates and FAO/OECD/IFPRI or national statistics may arise from different boundaries (e.g., what counts as “food” or “grain”) and from the fact that we estimate coefficients from survey-based and predicted ratios rather than from aggregate supply-use identities.

\textbf{Self-sufficiency and trade.} To the extent that total demand is higher than what at-home data alone suggest, self-sufficiency rates computed from at-home consumption will be overstated if the denominator ignores the out-of-home component. Our results do not directly yield a revised self-sufficiency rate, because that would require full supply-use balance and consistent treatment of stocks and trade; however, they indicate that any such balance should use total (at-home plus out-of-home) consumption, scaled by our coefficients, for a more accurate demand side. Our provincial-level outputs (\texttt{results\_province\_*.csv}) allow similar reasoning at the regional level for region-specific food security and trade analysis.

In sum, the paper's CSV-based estimates show that food away from home raises total consumption for all key food categories, with the largest relative effects for poultry, pork, legumes, and coarse grains. These findings are robust to the exclusion of imputation and are consistent with a growing role of dining out in Chinese diets. They support the use of total (at-home plus out-of-home) consumption in grain and food demand assessments and in discussions of China's food security and self-sufficiency.


% ========== 7. Conclusion ==========
\section{Conclusion}
\label{sec:conclusion}

This paper has provided the first machine-learning-based, commodity-level estimates of out-of-home food consumption coefficients for China at provincial and national level. We asked how much food China consumes away from home and operationalized the answer as the ratio of total to at-home consumption by 15 food categories and by province and year. Using micro household survey data and macro province-year covariates, we bridged the micro--macro gap with Copula income distribution matching, imputed missing consumption ratios using the best of 12 ML models (FT-Transformer), and applied spatial interpolation for provinces without micro data. We compared 12 models for coefficient prediction and reported results from LightGBM, TabPFN, and FT-Transformer as central estimates.

Main findings are threefold. First, out-of-home consumption coefficients are above one for all food categories and years considered, implying that food away from home adds non-negligibly to total demand. Second, the magnitude of the coefficient varies by category: poultry, pork, legumes, coarse grains, and aquatic products show larger multipliers; staples such as cooking oil and sugar show coefficients closer to one. Third, results are robust to excluding imputation: main and no-imputation (robust) specifications yield qualitatively similar conclusions, with point estimates differing by a few percentage points depending on model and category. The best imputation model in our selection is FT-Transformer (lowest MAE in cross-validation); TabPFN and FT-Transformer also perform well on the coefficient prediction task.

These results have direct implications for food security assessment. Because coefficients are above one, total grain and food demand is higher than at-home consumption alone. Policy and self-sufficiency analyses that ignore the out-of-home component will understate demand. Our provincial and national coefficient estimates can be used in food balance and demand projections and can be compared with FAO, OECD, IFPRI, and the literature where definitions and coverage align. Limitations include the following. First, survey coverage does not include all provinces or all years, so that some province--year cells rely entirely on spatial interpolation rather than on micro-based predictions; estimates for those cells are therefore more dependent on the interpolation method. Second, Copula matching assumes that the shape of the micro income distribution is transportable to macro cells (after scaling the mean), which may not hold if the underlying income distribution differs across regions or over time. Third, validation against independent macro-level out-of-home consumption data was not possible because such data are not available for China. Finally, different ML models yield somewhat different point estimates—though the bootstrap and robustness checks help quantify uncertainty and sensitivity.

% ========== 8. Policy Recommendations ==========
\section{Policy Implications and Recommendations}
\label{sec:policy}

Based on the evidence in this paper, we offer the following recommendations.

\textbf{Use total consumption in demand and food security assessments.} Government and international agencies that project grain and food demand or assess self-sufficiency should use total (at-home plus out-of-home) consumption, not at-home consumption alone. Our coefficients provide a consistent way to scale from at-home to total consumption by category and province-year. Incorporating these coefficients into existing food balance or demand models will improve the accuracy of demand-side estimates.

\textbf{Monitor provincial and category-level heterogeneity.} Out-of-home consumption varies by province and by food category. Policy aimed at food security or dietary change should take into account the provincial-level coefficients and per-capita quantities available in our \texttt{results\_province\_*.csv} outputs, so that regional targeting and resource allocation can be evidence-based.

\textbf{Align data collection with coefficient estimation.} Where possible, future household and macro surveys should record both at-home and total (or out-of-home) consumption by food category so that coefficient estimates can be updated and validated. Standardizing definitions of ``at home'' and ``away from home'' across surveys would facilitate comparison with our estimates and with international sources.

\textbf{Quantify uncertainty in policy scenarios.} Our bootstrap prediction intervals and robustness checks show that point estimates can vary by model and specification. When using our coefficients in scenario analysis, we recommend reporting results under more than one model (e.g., LightGBM and TabPFN or FT-Transformer) or using the bootstrap intervals where available, so that policy conclusions are not overly sensitive to a single point estimate.

\paragraph{Data and code availability.} The coefficient outputs (e.g., \texttt{results\_national\_*.csv}, \texttt{results\_province\_*.csv}, \texttt{model\_comparison\_national.csv}) and the code used to generate them are available in the supplementary material or from the authors upon request. Details on the underlying micro and macro data sources are provided in the Data subsection (Section~\ref{subsec:data}).

% ========== References ==========
\section*{References}
\addcontentsline{toc}{section}{References}

\bibliographystyle{apalike}
\bibliography{references}

\end{document}

% --- Approximate word count: body text (excluding abstract, tables, refs) ~8500--9000 words. ---
